\begin{helper}[Root component and nonisolation at a parallel step]
Let a certified Tashkinov series have a positive slot \(t\) of kind PE,
with chosen support \(v_t\), parallel color \(\gamma_t\), and defective
color \(\delta_t\).  Under the current coloring \(\varphi_t\), the two
colors are distinct and the coloring is proper.  The preceding tree
\(T_t\) contains \(v_t\) and is an initial segment of the next tree.
The \((\gamma_t,\delta_t)\)-component through \(v_t\) is a path
component whose only vertex in \(T_t\) is \(v_t\).  Moreover, for every
other vertex \(v\) of \(T_t\), the component through \(v\) contains an
exact edge identity.
\end{helper}
\begin{proof}
Inspect the certificate at slot \(t\) in \noderef{TashkinovSeries}.
The stated kind equation rules out the initial, SE, and RE constructors
of \noderef{TashkinovSeriesStep}, so the certificate is PE.  Its
chosen-vertex, gamma, and delta equations identify the PE constructor's
support, parallel color, and defective color with
\(v_t,\gamma_t,\delta_t\).  The constructor directly supplies
properness of \(\varphi_t\) and the required
\noderef{TreeSequenceInitialSegment} relation.

Write \(\varphi^-\) for the preceding coloring, \(\pi\) for the
selected coloring in \noderef{MaximumDefectiveChoice}, and \(\rho\) for
the second stable coloring.  The PE equations say
\[
  \varphi_t=\operatorname{swap}_{K_{v_t}(\gamma_t,\delta_t;\rho)}(\rho).
  \tag{1}
\]
We first show that \(V(T_t)\) is closed under \(\rho\).  It is closed
under \(\varphi^-\) by the PE certificate.  The selected coloring
\(\pi\) is proper, leaves exactly the distinguished edge uncolored, and
is stable from \(\varphi^-\), by
\noderef{MaximumDefectiveChoice}.  The preceding coloring has the same
two coloring properties.  Symmetry in \noderef{sc2} therefore gives
reverse stability.  By \noderef{BoundaryEdges}, every boundary edge is
incident with the tree.  If a boundary edge had \(\pi\)-color \(c\) missing
at a tree vertex under \(\pi\), the vertex clause of
\noderef{StableColoring} would make \(c\) missing at that vertex under
\(\varphi^-\), while reverse stability would fix that exact boundary
edge.  This contradicts the \noderef{ColoredVertexSetClosed} condition
for \(\varphi^-\).  Hence the tree is closed under \(\pi\).  Applying
the same argument to the PE's second stable relation, now between
\(\pi\) and \(\rho\), proves closedness under \(\rho\).

The selected edge \(f_t\) from \noderef{MaximumDefectiveChoice} is a
boundary edge incident with \(v_t\), has its other endpoint outside the
tree, and has \(\pi(f_t)=\delta_t\).  The second stable relation controls
\(\delta_t\), so \(\rho(f_t)=\delta_t\).  If \(\delta_t\) were missing
at any tree vertex \(x\), that same boundary identity would contradict
closedness under \(\rho\).  Thus \(\delta_t\) is present at every tree
vertex.  By the complementary definitions \noderef{MissingColors} and
\noderef{PresentColors}, for every such \(x\) there is an exact incident
edge \(g_x\) with \(\rho(g_x)=\delta_t\).

The \noderef{ParallelExtensionColor} clause says that \(\gamma_t\) is
missing at \(v_t\) under \(\pi\), and second stability transfers this
missingness to \(\rho\).  Since \(\delta_t\) is present there, the two
colors are distinct.  Unpack the PE's \noderef{TreeExit} witness.  Its
\noderef{TreeExitPath} is a simple two-color path from an outside vertex
to \(v_t\), puts \(v_t\) in \(T_t\), and makes \(v_t\) its only tree
vertex.  Missingness at both ends gives degree one there.  Every internal
node already has its two path neighbors, while
\noderef{TwoColorIncidenceDegreeAtMostTwo} bounds all degrees by two.  In the
graph of \noderef{TwoColorIncidenceGraph}, a shortest route from any further
node of the connected support \noderef{TwoColorComponent} to the path would
therefore create either a third neighbor at an internal node or a second
neighbor at an endpoint.  Hence the path exhausts the whole component
through \(v_t\).  It is consequently a
\noderef{TwoColorPathComponent}, with exactly the asserted tree
intersection.

The swap in (1), defined by \noderef{TwoColorComponentSwap}, changes
only this entire root component and merely exchanges its two colors.
By \noderef{TwoColorIncidenceGraph} and \noderef{TwoColorComponent}, exchanging
the two colors preserves the incidence graph and hence the component support.
Thus its unique tree vertex is unchanged under \(\varphi_t\).  For
\(x\ne v_t\), the edge-node of \(g_x\) cannot belong to the root component:
its incidence with \(x\) gives adjacency in \noderef{TwoColorIncidenceGraph},
which by the connected-support definition \noderef{TwoColorComponent} would
put \(x\) in that component,
contradicting the unique-intersection conclusion.  The swap therefore
fixes \(g_x\), so it remains a genuine \(\delta_t\)-colored exact edge
in the component through \(x\) under \(\varphi_t\).  This proves the
last assertion.
\end{proof}
